\expandafter\def\csname HistoricalNoteText9\endcsname{Ist $S>\BareRoot{\frac12}$, so folgt die zweite Bedingung aus der ersten; ist dagegen $S<\BareRoot{\frac12}$, so folgt die erste aus der zweiten.}

\expandafter\def\csname HistoricalNoteText10\endcsname{Es scheint gewiss hinreichend gut bewiesen zu sein, dass eine algebraische Kurve weder irgendwo plötzlich abbrechen kann (wie es etwa bei der transzendenten Kurve mit der Gleichung $y=\frac1{\mathrm{log}\,x}$ geschieht), noch sich nach unendlich vielen Windungen gleichsam in irgendeinem Punkt verlieren kann (wie die logarithmische Spirale); und meines Wissens hat niemand dagegen einen Zweifel erhoben. Sollte dennoch jemand einen Beweis verlangen, der keinerlei Zweifeln ausgesetzt ist, so werde ich es übernehmen, ihn bei einer anderen Gelegenheit zu geben. Im vorliegenden Fall ist aber klar: Verließe irgendein Zweig, etwa 2, den Kreis nirgends wieder (Fig. 3), so könnte man zwischen 0 und 2 in den Kreis eintreten, dann um diesen ganzen Zweig herumgehen (der sich im Kreisgebiet verlieren müsste) und schließlich zwischen 2 und 4 wieder aus dem Kreis austreten, ohne auf dem ganzen Weg irgendwo die erste Linie zu treffen. Dass dies widersinnig ist, erkennt man daran, dass man beim Eintritt in den Kreis die erste Fläche über sich hatte, beim Austritt aber unter sich. Man müsste daher notwendig irgendwo auf die erste Fläche selbst, also auf einen Punkt der ersten Linie getroffen sein. --- Aus dieser Überlegung, die sich auf Grundsätze der Geometrie der Lage stützt, die nicht weniger gültig sind als die der Geometrie der Größe, folgt übrigens nur, dass man, wenn man auf irgendeinem Zweig der ersten Linie in den Kreis eintritt, an einer anderen Stelle wieder austreten kann, wobei man stets auf der ersten Linie bleibt; nicht aber, dass der Weg eine stetige Linie in dem Sinne ist, wie man es in der höheren Geometrie versteht. Hier genügt jedoch, dass der Weg im gewöhnlichen Sinne eine stetige Linie ist, d.\,h. nirgends unterbrochen, sondern überall zusammenhängend.}
